\documentclass{article}
\usepackage[a4paper, total={6in, 10in}]{geometry}

\usepackage{lipsum}

% \usepackage{float}
\usepackage[]{graphicx}
\graphicspath{{Images/}}
\usepackage{hyperref}
% \usepackage{subfiles}

\usepackage{biblatex}
\addbibresource{ref.bib}

% \usepackage{booktabs}
% \usepackage{multicol}

% \usepackage{tabularx}

% \setlength\parindent{0pt}

% remove section numbering
% \setcounter{secnumdepth}{0}

\begin{document}
\pagenumbering{arabic}


\title{Working Document of Project}
\author{Ali Raeisdanaei, Jingyue Zhang, Tiantian Lin, Ziqian Qiu }
\date{}
\maketitle

\section{Preface}
This is a drafting document created for the purpose of writing any ideas down.
This document is only a draft; therefore, it is rough and messy.
However, the full report should evolve out of this document.


\section{Introduction}

\subsection{Purpose of this}
This project seeks to lay a ground build by systematic literature review for the layman and the practitioner to understand the free and open source hardware (FOSH).
It assumes no knowledge of the free and open source movement, nor hardware.
Through explaining all of this and more from the available literature, the project seeks to answer the following \textbf{Research Questions}:
RQ1: What is the state-of-the art FOSH and what are the performance of this best FOSH? How does it compare in terms of performance and other aspects with the non-FOSH state-of-the-art? 

RQ2: What are the main challenges and drawbacks associated with the development, adoption, and sustainability of FOSH?

RQ3: How does non-FOSH compare with FOSH in terms of design and licsensing?

RQ4: To what extent is FOSS more adapted FOSH development, and what are the factors that contribute to any observed differences in their adaptability?

RQ5: What are the potential future developments, opportunities, and challenges that FOSH is facing, and what are the implications of these for the growth and sustainability of the FOSH movement?

\subsection{Where is the Freedom?}
We have on our hands a revolution of technology. 
And it is not the endless march of \textit{forward} preached by exploitive companies.
This revolution is about \textbf{our} freedom to use technology for our learning, fun, and betterment.
We use technology; technology does not use us. 
Freedom and liberty over our technology is the extension of our freedom and liberty over our personhood.
This revolution is championed by the Richard Stallman, computer scientist, philosopher, preacher, saint, and the founder of GNU and the Free and Open Source Software organisation \cite{b0_stallman}.
As the \textit{popular} song goes \textit{Remember Richard Stallman,
Who set your software free} \cite{song_GNUs_Not_Unix}.

\begin{figure}[h]
    \centering
    \includegraphics[width=\textwidth]{the_wise_gnu.jpg}
    \caption{Richard Matthew Stallman as depicted in lore}
    \label{fig:rms}
\end{figure}

The passion and revolution of this movement, for users to take back ownership of their technology, is incredibly powerful. 
No literature on the topic of open source anything can distance itself with this revolutionary zeal. 
While it has not been fully actualised, it has permeated many aspects of our lives over the years.
We have the rise of platforms such as GitHub and the widespread adoption of the Public Licenses preached by the Free and Open Source movement. 
Examples of its victories are the Linux kernel, Mozilla Firefox, \LaTeX, Vim, and other widespread software. 

To understand the definition of the \textit{free and open source} in free and open source hardware, one must turn to its founding essays. Here is a good working definition of this:

TODO actually write the definition \cite{b0_stallman}.

Open source is not the same thing as free and open source [CITATION REQUIRED], and to not do injustice to the ideals of the revolution, this project, unlike much of the literature, will not drop the \textit{'F'} in FOSH nor in writing nor in definition.
This theory and youthful, revolutionary spirit will guide the project as an anchor.

\subsection{Where is the Hardware?}
An overview of hardware and its difficulties is needed as a good background information. 

The original Free and Open Source movement was about software not hardware.
% There needs to then be an extension of the principles to Hardware.


\section{TODOs}
% We need an introduction to what free and open source is. 
% A history of this in the software realm and how this idea is extending to hardware is needed. 

% We need to understand what are the values and principles of free and open source, and how they exactly relate to free and open source hardware (FOSH).
% To this end we should consult the founder of it all Richard Stallman (RMS).
% This should be a good source to read into
% \textbf{Free Software, Free Society: Selected Essays of Richard M. Stallman}
% \cite{b0_stallman}.


Then, we also need to explain why a literature review is needed.
What exactly does our work add?

\section{Readings done by Ali}

I had vastly underestimated the amount of work in this field, and now I understand why a systematic literature review is needed.

\subsection{Defining success in open source hardware development projects: a survey of practitioners}
\cite{p1_def_succ}

\textbf{Very Important Paper for all of us to Read}

This paper lays out some very crucial starting point information.
It even has a section of literature review for us to start.

\subsubsection{Open Source Hardware Association}
There is an open source hardware association that tracks these hardware.
\href{https://www.oshwa.org}{Open Source Hardware Association}
This is a very crucial resource for us to dive into. 
The association gives out certification for projects that meet its requirements
\href{https://certification.oshwa.org/list.html}{https://certification.oshwa.org/list.html}

\subsubsection{Open Source Hardware Journals}
There have also been journals that have been dedicated to the FOSH.
\begin{enumerate}
    \item \href{https://openhardware.metajnl.com/}{Journal of Open Hardware}
    \item \href{https://www.sciencedirect.com/journal/hardwarex}{HardwareX}
\end{enumerate}

\subsubsection{Summary of Paper}

\section{Methodology of Literature Review}
We first identified a seed of papers, and recursively read the papers that were cited until saturation was reached. 

\nocite{*}
\printbibliography

\end{document}